% Source: https://artofproblemsolving.com/community/c3h2226421_method_for_solving_modulo_congruences_likely_flawed
% Stable identity: aops-topic-2226421-q01
A student proposes this method: the solutions of the system
\[x\equiv n\pmod y,\qquad x\equiv n\pmod z\]
are $x=n+t\cdot\operatorname{lcm}(y,z)$ for positive integers $t$. For example, $x\equiv10\pmod{11}$ and $x\equiv10\pmod{13}$ give $x=10+143t$.
\begin{enumerate}
\item[(a)] Prove that the integers $x$ satisfying the system are exactly the $x$ with $x\equiv n\pmod{\operatorname{lcm}(y,z)}$, and that the same holds for any number of congruences with the same remainder $n$.
\item[(b)] What is wrong with ``for positive integers $t$''?
\end{enumerate}
